\documentclass{beamer}
\usepackage{etex}
\usepackage{lmodern}
\usepackage[english]{babel}
\usepackage[T1]{fontenc}
\usepackage[utf8]{inputenc}
\usepackage{pst-sigsys} 
\usepackage{amsmath,amsfonts,amssymb}
\usepackage{pstricks-add} 
\usepackage{ragged2e}
\usepackage{graphicx}
\usepackage{ulem}
\usepackage{fontawesome}
\setbeamertemplate{navigation symbols}{}  
 
\usetheme{Darmstadt}
\setbeamertemplate{footline}{\insertframenumber/\inserttotalframenumber}
\title{M1 IEAP - BTI/FH/IEMH\\pFIEA02CM : Analyse et Traitement du Signal}
\author{Flavy ROSEREN\\Martin EGIZIANO\\Frank BULOUP}
\institute{Aix Marseille Université\\Institut des Sciences du Mouvement}
\date{}

\setbeamertemplate{footline} 
{  
	\begin{beamercolorbox}[ht=2.5ex,dp=1.125ex,%
      leftskip=.3cm,rightskip=.3cm plus1fil]{title in head/foot}%
      {\usebeamerfont{title in head/foot}\insertshorttitle} \hfill    
      \insertframenumber / \inserttotalframenumber%
    \end{beamercolorbox}%
%     \begin{beamercolorbox}[colsep=1.5pt]{lower separation line foot}
%     \end{beamercolorbox} 
}

\newcounter{exampleBlockCounter}
\setcounter{exampleBlockCounter}{1}

\begin{document} 
 
\begin{frame}[plain] 
	\titlepage 
	\center{\includegraphics[scale=0.75]{images/by-nc-sa.eps}}
	\vspace{1cm}
	\includegraphics[scale=0.6]{images/LogoAMU.png}\hspace*{2cm}
	\includegraphics[scale=0.2]{images/LogoCNRS.eps}\hspace*{2cm}
	\includegraphics[scale=0.1]{images/LogoISM.eps}
\end{frame}

\begin{frame}{Première Partie}
	\tableofcontents 
\end{frame}

\section{Signal discret - Aspect fréquentiel}
\subsection{Série de Fourier} 
\begin{frame}
\begin{block}{Généralisation}
\justifying L'exemple de l'exercice précédent peut se généraliser  :
\textcolor{red}{presque tout \textbf{signal continu} périodique peut être
synthétisé par sommation de sinus et/ou de cosinus}. C'est ce que l'on appelle la \textcolor{red}{Série de Fourier}. Un signal $s(t)$ de
période $T_0 = \frac{1}{f_0}$ peut alors s'écrire de la façon suivante :
\begin{align*}
s(t) = S(0) + \sum_{k=1}^{+\infty} 2\rho (k) cos[2\pi kf_0t+\phi (k)]
\end{align*}
\end{block}
\pause
\begin{alertblock}{Remarques}
	\itemize
	\item $\rho(k)$ et $\phi(k)$ sont respectivement des amplitudes et des phases
	\pause
	\item $S(0)$ est la valeur moyenne, c'est toujours un réel pur
	\pause
	\item $f_0$ est la fréquence fondamentale 
	\pause
	\item Les fréquences des cosinus sont des multiples entiers de $f_0$. Ce sont
	les harmoniques.
\end{alertblock}
\end{frame}

\begin{frame}
\begin{alertblock}{Mais comment obtenir $\rho(k)$ et $\phi(k)$ ?}
\justifying Par le calcul du nombre complexe $S(k)$ suivant :
\[
S(k) = \frac{1}{T_0}\int_{-T_{0}/2}^{T_{0}/2}s(t)e^{-2i\pi k f_{0}t} dt
\Rightarrow
\begin{cases}
\rho(k) = \lvert S(k) \rvert\\
\phi(k) = Arg[S(k)]
\end{cases}
\]
% \justifying Qui donnera $\rho(k) = \lvert S(k) \rvert$ et $\phi(k) = Arg[S(k)]$
\end{alertblock}
\pause
\begin{block}{Série de Fourier du signal carré}
\justifying Le calcul de $S(k)$ sur le signal carré précédent nous donnerait :
\[
S(k) =
\begin{cases}
 \frac{2}{k\pi}e^{-i\frac{\pi}{2}}\text{ pour toutes les valeurs de
}k\text{ impaires}\\
0\text{ pour toutes les valeurs de }k\text{ paires}
\end{cases}
\]
\end{block}
\end{frame} 

\begin{frame}
\begin{alertblock}{Remarques}
\itemize \justifying
\item La synthèse d'un signal périodique par une série de Fourier
nécessite un nombre infini d'harmoniques, c.a.d. que
le signal est composé de bien plus hautes fréquences que la
fondamentale.
\pause
\item En pratique on se limite a un nombre fini de termes. Le signal
obtenu est alors une approximation du signal périodique initial.
\pause
\item Si le signal initial présente des points de discontinuités, un phénomène
oscillatoire autour de ces points vient compromettre la convergence : en ces
points, la série de Fourier ne converge pas vers le signal original. C'est ce
que l'on nomme \textbf{phénomène de Gibbs}.
\end{alertblock}
\end{frame}

\subsection{Représentation fréquentielle des signaux}
\begin{frame}
\begin{block}{Représentation fréquentielle et spectres monolatéraux}
\justifying
Le développement d'un signal en série de Fourier donne donc une information
sur sa composition en fréquence : le signal est une somme de cosinusoïdes dont
les amplitudes, les phases et les fréquences sont bien définies.
Il est donc possible de le caractériser par un ensemble de triplets $[2\rho(k),
\phi(k), kf_0]$ lorsque $k$ est différent de zéro et par le triplet $[\rho(0),
\phi(0), 0]$ pour la fréquence nulle, $\phi(0)$ valant $0$ ou $\pi$.
\end{block}
\pause
\begin{columns}[t]
	\begin{column}{4cm}
		\begin{figure}
			\begin{pspicture}[showgrid=false](0,-.5)(4,3)
				\psaxeslabels{->}(0,0)(0,0)(4,3){$f$}{$Amplitude$}
				\psstem[style=Stem]{2,.5,1,0}
				\rput(0,-.250){$0$}
				\rput(1,-.250){$f_0$}
				\rput(2,-.250){$2f_0$}
				\rput(3,-.250){$3f_0$}
				\rput(.450,2.25){$\rho(0)$}
				\rput(1,.850){$2\rho(1)$}
				\rput(2,1.35){$2\rho(2)$}
				\rput(3,0.35){$2\rho(3)$}
				\pause
				\rput*{45}(2,1.75){\textbf{\textcolor{red}{Spectre d'amplitude}}}
			\end{pspicture}
		\end{figure}
	\end{column}
				\pause
	\begin{column}{4cm}
		\begin{figure}
			\begin{pspicture}[showgrid=false](0,-2.5)(4,1)
				\psaxeslabels{->}(0,0)(0,-2)(4,1){$f$}{$Phase$}
				\psstem[style=Stem]{0,-1.5,-.25,0}
				\rput(-.250,.250){$0$}
				\rput(1,.250){$f_0$}
				\rput(2,.250){$2f_0$}
				\rput(3,.250){$3f_0$}
				\rput(.40,-.25){$\phi(0)$}
				\rput(1,-1.850){$\phi(1)$}
				\rput(2,-.65){$\phi(2)$}
				\rput(3,-0.35){$\phi(3)$}
				\pause
				\rput*{45}(2,-.75){\textbf{\textcolor{red}{Spectre de phase}}}
			\end{pspicture}
		\end{figure}
	\end{column}
\end{columns}
\end{frame}

\begin{frame}
\begin{exampleblock}{Exercice \Roman{exampleBlockCounter} - Représentation spectrale monolatérale}
Représenter les spectres monolatéraux des signaux suivants :
\begin{itemize}
  \item $s_1(t) = 3$
  \item $s_2(t) = 6sin(2\pi t)$
  \item $s_3(t) = 1 + 2cos(10\pi t) + 4cos(20\pi t)$
  \item $s_4(t) = cos(100\pi t) + cos^2(100\pi t)$
\end{itemize}
\end{exampleblock}
\end{frame}
\stepcounter{exampleBlockCounter}
\begin{frame}
\begin{exampleblock}{Exercice \Roman{exampleBlockCounter} - Représentation spectrale monolatérale}
\itemize\justifying
\item Représenter les spectres monolatéraux d'amplitude et de phase du signal carré
\item Écrire le signal temporel associé aux spectres monolatéraux :
\end{exampleblock}
\begin{columns}[t]
	\begin{column}{4cm}
		\begin{figure}
			\begin{pspicture}[showgrid=false](0,-.5)(4,3)
				\psaxeslabels{->}(0,0)(0,0)(4,3){$f$}{$Amplitude$}
				\psstem[style=Stem]{2,.5,1,0}
				\rput(0,-.250){$0$}
				\rput(1,-.250){$f_0$}
				\rput(2,-.250){$2f_0$}
				\rput(3,-.250){$3f_0$}
				\rput(.450,2.25){$4$}
				\rput(1,.850){$1$}
				\rput(2,1.35){$2$}
				\rput(3,0.35){$ $}
			\end{pspicture}
		\end{figure}
	\end{column}
	\begin{column}{4cm}
		\begin{figure}
			\begin{pspicture}[showgrid=false](0,-2.5)(4,1)
				\psaxeslabels{->}(0,0)(0,-2)(4,1){$f$}{$Phase$}
				\psstem[style=Stem]{0,-1.5,-.25,0}
				\rput(-.250,.250){$0$}
				\rput(1,.250){$f_0$}
				\rput(2,.250){$2f_0$}
				\rput(3,.250){$3f_0$}
				\rput(.40,-.25){$ $}
				\rput(1,-1.850){$-3$}
				\rput(2,-.65){$-0.5$}
				\rput(3,-0.35){$ $}
			\end{pspicture}
		\end{figure}
	\end{column}
\end{columns}
\end{frame}

\stepcounter{exampleBlockCounter}

\begin{frame}
\begin{exampleblock}{Exercice \Roman{exampleBlockCounter} - Série de Fourier}
\justifying Retrouver le résultat de l'exercice précédent par le calcul de
$S(k)$ sur le signal suivant :
\begin{figure}
% 	\psset{xunit=1cm}
% 	\psset{yunit=1cm}
	\begin{pspicture}[showgrid=false](-5,-1)(5,1.25)
		\psaxeslabels[Dx=1]{->}(0,0)(-5,-1)(5,1.25){$t$}{$s(t)$}
		\psline(0,0.5)(1,0.5)
		\psline(1,-0.5)(2,-0.5)
		\psline(2,0.5)(3,0.5)
		\psline(3,-0.5)(4,-0.5)
		\psline(-1,-0.5)(0,-0.5)
		\psline(-2,0.5)(-1,0.5)
		\psline(-3,-0.5)(-2,-0.5)
		\psline(-4,0.5)(-3,0.5)
		\psline(1,-0.5)(1,0.5)
		\psline(2,-0.5)(2,0.5)
		\psline(3,-0.5)(3,0.5)
		\psline(-1,-0.5)(-1,0.5)
		\psline(-2,-0.5)(-2,0.5)
		\psline(-3,-0.5)(-3,0.5)
		\rput(4.5,0.25){\Large\ldots}
		\rput(-4.5,0.25){\Large\ldots}
		\rput(2.25,-0.25){$T_0$}
		\rput(-.25,0.5){$1$}
		\rput(0.25,-0.5){$-1$}
	\end{pspicture}
\end{figure}
\end{exampleblock}
\end{frame}

\end{document}
